% Introduction
\section{Introduction}
\label{sec:intro}

Multi-Agent Path Finding (MAPF) asks for a set of collision-free paths for a group of agents navigating from their respective start configurations to goal configurations~\cite{stern2019mapfsurvey}. Classical MAPF operates on graphs, where agents occupy discrete nodes and move along edges in synchronized time steps. This abstraction enables powerful algorithmic solutions such as Conflict-Based Search (CBS)~\cite{sharon2015cbs} and Priority Inheritance with Backtracking (PIBT)~\cite{okumura2022pibt}, but it imposes a fundamental mismatch with physical deployments: real robots move through continuous space with bounded velocities, non-trivial body radii, and smooth kinematic constraints that discrete graph abstractions cannot faithfully represent.

Continuous-space MAPF requires every agent to simultaneously (i) generate a kinematically feasible trajectory toward its goal and (ii) ensure that no two agents occupy overlapping regions at any point in time. Classical continuous approaches such as Optimal Reciprocal Collision Avoidance (ORCA)~\cite{vandenberg2011orca} and its predecessors~\cite{vandenberg2008rvo} handle the collision avoidance problem via local velocity-space constraints, but they optimize myopically, ignoring global path structure. Under high agent density, ORCA frequently deadlocks because no single agent can satisfy all its velocity constraints simultaneously. Learning-based methods that predict preferred velocities can improve global coordination~\cite{sartoretti2019primal,ma2021dhc,wang2023scrimp}, but without a hard safety layer, they accumulate collisions especially in dense or unseen environments.

This paper addresses the gap between expressive learned motion planning and rigorous conflict resolution in continuous space. We make two primary contributions. First, we introduce \textbf{EPIBTShield}, a continuous-space generalization of the discrete PIBT algorithm~\cite{okumura2022pibt}. EPIBTShield maintains a dynamic priority ordering over agents and resolves velocity conflicts by selecting from a scored discrete set of candidate velocities in strict priority order, with recursive backtracking to avoid global deadlock. Unlike ORCA, EPIBTShield does not require the velocity obstacle set to be non-empty; it always produces a valid (possibly zero) velocity for every agent while respecting hard geometric collision constraints. Second, we train a \textbf{rectified-flow GNN} that maps each agent's egocentric local observation to a preferred velocity. The network is conditioned on a graph of neighboring agents and is trained via flow matching~\cite{liu2023rectifiedflow,lipman2023flowmatching} to imitate expert trajectories generated by the EECBS planner~\cite{li2021eecbs}. These preferred velocities serve as inputs to EPIBTShield, providing goal-directed, spatially-aware guidance that substantially reduces the reliance on the shield's fallback candidates.


Our empirical evaluation spans four map types and two agent counts drawn from the MovingAI benchmark suite~\cite{sturtevant2012movingai}. On the primary in-distribution benchmark (random-32-32-10, $N=100$), the combined \mbox{Flow+EPIBTShield} system achieves an AtGoal fraction of $0.874$, compared to $0.168$ for ORCA and $0.486$ for a flow model shielded by ORCA — a gap of $70.6$ percentage points over the strongest pure-ORCA baseline. We also report honestly on the system's limitations: in warehouse-style out-of-distribution maps, the learned flow component underperforms a simple straight-line preferred velocity, suggesting that the generalization of the flow model is bounded by the distribution of its training demonstrations.

\subsection*{Summary of Contributions}

\begin{enumerate}
  \item \textbf{EPIBTShield}: a continuous-space velocity-conflict resolution algorithm that extends PIBT~\cite{okumura2022pibt} to continuous domains via dynamic priority ordering, a cosine-aligned candidate-velocity set with $19$ candidates per agent, and recursive backtracking up to depth~3.

  \item \textbf{Rectified-flow GNN}: a six-layer GraphSAGE~\cite{hamilton2017graphsage} network with egocentric CNN feature extraction, trained with rectified flow~\cite{liu2023rectifiedflow} on EECBS~\cite{li2021eecbs} expert demonstrations, that predicts spatially-informed preferred velocities for EPIBTShield.

  \item \textbf{Empirical evaluation} on MovingAI benchmarks~\cite{sturtevant2012movingai} demonstrating $70.6$\,pp improvement over ORCA in dense random maps, with systematic ablations and an honest analysis of OOD failure modes and multi-seed training variance.
\end{enumerate}

The paper is organized as follows. Section~\ref{sec:related} reviews related work. Section~\ref{sec:method} describes EPIBTShield and the flow GNN in detail. Section~\ref{sec:experiments} presents quantitative results and ablations. Section~\ref{sec:conclusion} concludes.
